% GENERATED FILE. Edit canonical lessons, then run npm run generate:books.
% canonical-source-hash: fnv1a64:d92787afe4b795b7
% canonical-lessons: FR-C501-collectionner, FR-C501-collectionneur, FR-C501-maquette, FR-C501-modelisme, FR-C501-assembler

\chapter{Collectionneurs}
\label{ch:fr-collectionneurs}
\hlchaptermodality{501}

\begin{chapteropening}
\textbf{By the end of this chapter:} \emph{I can talk about collecting and building models.}

Complete the last lesson of chapter 501: I can talk about collecting and building models.
\end{chapteropening}

\section[collectionner]{collectionner — to collect}

\label{lesson:FR-C501-collectionner}



\begin{quote}
Before the new one: say the French for bet, then the French for chance.
\end{quote}

\subsection*{You'll want to know: collectionner}
\textbf{collectionner} — \textquotedblleft{}to collect\textquotedblright{}.

\begin{quote}
Mon voisin collectionne les vieilles cartes postales.
\end{quote}

\begin{quote}
\emph{My neighbour collects old postcards.}
\end{quote}

As a hobby. To pick up someone is \emph{aller chercher}.

\begin{cousinweb}
\emph{collectionner} is built on \emph{collection} (EtymDB), a learned borrowing from Latin \emph{collēctiō} \textquotedblleft{}a gathering\textquotedblright{}, from \emph{colligere} \textquotedblleft{}to gather\textquotedblright{}, built on \emph{legere} \textquotedblleft{}to gather, to read\textquotedblright{} (Lewis \& Short).
\end{cousinweb}

\begin{grammarlens}[title={What you've built}]
One more word for telling, explaining and describing at B1.
\end{grammarlens}

\subsection*{Your turn}
\begin{itemize}
\raggedright
  \item \emph{Say it:} \emph{collectionner}
  \item \emph{Say it:} \emph{collectionner}, once more
  \item \emph{Say it:} a sentence of your own with \emph{collectionner}
  \item \emph{Recall:} say the French for bet, then the French for chance, then say \emph{collectionner} again
\end{itemize}

\begin{tcolorbox}[breakable,colback=teal!4,colframe=teal!35!black,title={Before you move on}]
What does collectionner mean? (\textquotedblleft{}To collect\textquotedblright{}.) Say it once more.
\end{tcolorbox}
\section[le collectionneur]{le collectionneur — collector}

\label{lesson:FR-C501-collectionneur}



\begin{quote}
Before the new one: say the French for chance, then the French for to collect.
\end{quote}

\subsection*{You'll want to know: le collectionneur}
\textbf{le collectionneur} — \textquotedblleft{}collector\textquotedblright{}.

\begin{quote}
Ce collectionneur possède deux mille timbres.
\end{quote}

\begin{quote}
\emph{This collector owns two thousand stamps.}
\end{quote}

A woman is \emph{la collectionneuse}.

\begin{cousinweb}
\emph{le collectionneur} is built on \emph{collectionner} \textquotedblleft{}to collect\textquotedblright{}, from \emph{collection}, a learned borrowing from Latin \emph{collēctiō}, built on \emph{legere} \textquotedblleft{}to gather\textquotedblright{} (Lewis \& Short).
\end{cousinweb}

\begin{grammarlens}[title={What you've built}]
One more word for telling, explaining and describing at B1.
\end{grammarlens}

\subsection*{Your turn}
\begin{itemize}
\raggedright
  \item \emph{Say it:} \emph{le collectionneur}
  \item \emph{Say it:} \emph{le collectionneur}, once more
  \item \emph{Say it:} a sentence of your own with \emph{le collectionneur}
  \item \emph{Recall:} say the French for chance, then the French for to collect, then say \emph{le collectionneur} again
\end{itemize}

\begin{tcolorbox}[breakable,colback=teal!4,colframe=teal!35!black,title={Before you move on}]
What does le collectionneur mean? (\textquotedblleft{}Collector\textquotedblright{}.) Say it once more.
\end{tcolorbox}
\section[la maquette]{la maquette — scale model; mock-up}

\label{lesson:FR-C501-maquette}



\begin{quote}
Before the new one: say the French for to collect, then the French for collector.
\end{quote}

\subsection*{You'll want to know: la maquette}
\textbf{la maquette} — \textquotedblleft{}scale model; mock-up\textquotedblright{}.

\begin{quote}
Il a construit une maquette de bateau.
\end{quote}

\begin{quote}
\emph{He built a model boat.}
\end{quote}

Also a designer's mock-up of a page or a building.

\begin{cousinweb}
\emph{la maquette} was borrowed from Italian \emph{macchietta} \textquotedblleft{}a little sketch\textquotedblright{} (EtymDB), from \emph{macchia} \textquotedblleft{}a spot\textquotedblright{}, inherited from Latin \emph{macula} \textquotedblleft{}a spot\textquotedblright{} (Lewis \& Short).
\end{cousinweb}

\begin{grammarlens}[title={What you've built}]
One more word for telling, explaining and describing at B1.
\end{grammarlens}

\subsection*{Your turn}
\begin{itemize}
\raggedright
  \item \emph{Say it:} \emph{la maquette}
  \item \emph{Say it:} \emph{la maquette}, once more
  \item \emph{Say it:} a sentence of your own with \emph{la maquette}
  \item \emph{Recall:} say the French for to collect, then the French for collector, then say \emph{la maquette} again
\end{itemize}

\begin{tcolorbox}[breakable,colback=teal!4,colframe=teal!35!black,title={Before you move on}]
What does la maquette mean? (\textquotedblleft{}Scale model; mock-up\textquotedblright{}.) Say it once more.
\end{tcolorbox}
\section[le modélisme]{le modélisme — model-making}

\label{lesson:FR-C501-modelisme}



\begin{quote}
Before the new one: say the French for collector, then the French for scale model.
\end{quote}

\subsection*{You'll want to know: le modélisme}
\textbf{le modélisme} — \textquotedblleft{}model-making\textquotedblright{}.

\begin{quote}
Mon père fait du modélisme : il construit des petits trains.
\end{quote}

\begin{quote}
\emph{My father does model-making: he builds little trains.}
\end{quote}

\emph{Un modéliste} — a model-maker.

\begin{cousinweb}
\emph{le modélisme} is built on \emph{modèle} \textquotedblleft{}model\textquotedblright{} with \emph{-isme}; \emph{modèle} was borrowed from Italian \emph{modello}, which goes back to Latin \emph{modulus} \textquotedblleft{}a small measure\textquotedblright{}, from \emph{modus} \textquotedblleft{}measure\textquotedblright{} (Lewis \& Short).
\end{cousinweb}

\begin{grammarlens}[title={What you've built}]
One more word for telling, explaining and describing at B1.
\end{grammarlens}

\subsection*{Your turn}
\begin{itemize}
\raggedright
  \item \emph{Say it:} \emph{le modélisme}
  \item \emph{Say it:} \emph{le modélisme}, once more
  \item \emph{Say it:} a sentence of your own with \emph{le modélisme}
  \item \emph{Recall:} say the French for collector, then the French for scale model, then say \emph{le modélisme} again
\end{itemize}

\begin{tcolorbox}[breakable,colback=teal!4,colframe=teal!35!black,title={Before you move on}]
What does le modélisme mean? (\textquotedblleft{}Model-making\textquotedblright{}.) Say it once more.
\end{tcolorbox}
\section[assembler]{assembler — to put together, to assemble}

\label{lesson:FR-C501-assembler}



\begin{quote}
Before the new one: say the French for scale model, then the French for model-making.
\end{quote}

\subsection*{You'll want to know: assembler}
\textbf{assembler} — \textquotedblleft{}to put together, to assemble\textquotedblright{}.

\begin{quote}
Il a assemblé l'armoire en deux heures.
\end{quote}

\begin{quote}
\emph{He put the wardrobe together in two hours.}
\end{quote}

\emph{S'assembler} — to gather.

\begin{cousinweb}
\emph{assembler} is inherited from Latin \emph{assimulāre} (EtymDB), which in Classical Latin meant \textquotedblleft{}to make like\textquotedblright{}; the sense \textquotedblleft{}to bring together\textquotedblright{} is usually traced to \emph{simul} \textquotedblleft{}together\textquotedblright{} (Lewis \& Short).
\end{cousinweb}

\begin{grammarlens}[title={What you've built}]
One more word for telling, explaining and describing at B1.
\end{grammarlens}

\subsection*{Your turn}
\begin{itemize}
\raggedright
  \item \emph{Say it:} \emph{assembler}
  \item \emph{Say it:} \emph{assembler}, once more
  \item \emph{Say it:} a sentence of your own with \emph{assembler}
  \item \emph{Recall:} say the French for scale model, then the French for model-making, then say \emph{assembler} again
\end{itemize}

\begin{tcolorbox}[breakable,colback=teal!4,colframe=teal!35!black,title={Before you move on}]
What does assembler mean? (\textquotedblleft{}To put together, to assemble\textquotedblright{}.) Say it once more.
\end{tcolorbox}
